\subsection{Decorated witness}

\begin{theorem}[Jeandel--Rao finite aperiodic cohomology witness]\label{thm:jeandel-rao-witness}\lean{Obsmetrologia.AperiodicOrder.Tiling.Witness.jeandelRaoMarker01_nonzero_h1_class}\leanok\uses{thm:clinv-dichotomy}
A concrete, finite realization of ``$\mathrm{cl\text{-}inv}$ can be nonzero in
the aperiodic regime'': on a finite one-vertex Anderson--Putnam
marker-domino subcomplex extracted from the Jeandel--Rao $11$-tile aperiodic
Wang set, the decorated cochain $\varphi$ supported on $\{e_2,e_3\}$ is a
nonzero cocycle. Since the one-vertex complex has $\delta^0 = 0$,
$[\varphi]\neq0$ in $H^1$.

This is a finite witness only: it is not Jeandel--Rao aperiodicity, and not by
itself a persistent direct-limit $\mathrm{cl\text{-}inv}$ class. It is,
however, the cohomological content of the abstract
\texttt{aperiodicWitness}: \texttt{Settling.clInv} is derived from a closed
cochain on a one-vertex AP complex, and
\texttt{aperiodicWitness\_clinv\_nonzero} is discharged from
\texttt{phi\_nonzero} rather than from a bare coordinate.
\end{theorem}
